\section{ELEMENTS AND SUBSETS OF A SET}

\begin{enumerate}
\def\labelenumi{\arabic{enumi}.}
\item
  A set consists of elements which are capable of possessing certain properties and of having certain relations between themselves or with elements of other sets.
\item
  Sets and elements are denoted in mathematical arguments by graphical symbols, which in general are letters (from various alphabets) or combinations of letters and other signs. Relations between elements of one or more sets are denoted by inserting the symbols which denote these elements into a scheme characteristic of the relation considered (†); and similarly for properties.
\end{enumerate}

A letter may denote either a fixed element or an arbitrary element (also called a variable, an argument, or a generic element) of a set. When an arbitrary element is replaced by a fixed element in a relation (or property), the arbitrary element is said to be given this fixed element as value.

In order to indicate the elements which appear in a relation which is not explicitly written down, we represent the relation by a notation such as \(R\{x,y,z\}\) (if x,y,z are the elements in question).

\begin{enumerate}
\def\labelenumi{\arabic{enumi}.}
\setcounter{enumi}{2}
\item
  A relation or a property in which arbitrary elements feature (*) is said to be an identity if it becomes a true proposition whatever values we give to these arbitrary elements. If R and S denote two relations (or properties), R is said to imply S if S is true whenever the arbitrary elements which enter in these relations are fixed in such a way that R is true. The relations (or properties) R and S are said to be equivalent if each implies the other.
\item
  Let \(R\{x, y, z\}\) be a relation between the variables x, y, z. The phrase ``for all x, \(R\{x, y, z\}\) '' is a relation between y and z, which will be considered to be true for a system of given values of these latter variables if R is true for these values of y and z and every value of x. The phrase ``there exists x such that \(R\{x, y, z\}\) '' (or ``for some x, \(R\{x, y, z\}\) '') is again a relation between y and z, which will be considered to be true for a system of given values of y and z if, these variables being thus fixed, there is at least one value of x for which R is true. Similarly for a relation between any number of variables.
\end{enumerate}

If \(\overline{R}\) denotes the negation of R, then the negation of ``for all x, R'' is ``there exists x such that \(\overline{R}\) ''; the negation of ``there exists x such that R'' is ``for all x, \(\overline{R}\) ''.

\begin{enumerate}
\def\labelenumi{\arabic{enumi}.}
\setcounter{enumi}{4}
\item
  If R, S denote two relations, we regard ``R and S'' as a single relation, which is considered as true whenever both R and S are true. Likewise, ``R or S'' is a relation which is considered to be true whenever at least one of the relations R, S is true (and, in particular, whenever they are both true). The word ``or'' thus does not have the disjunctive sense here which it sometimes has in ordinary speech. Let \(\overline{R}\) , \(\overline{S}\) denote the negations of R, S, respectively. Then the negation of ``R and S'' is `` \(\overline{R}\) or \(\overline{S}\) '', and the negation of ``R or S'' is `` \(\overline{R}\) and \(\overline{S}\) ''.
\item
  By writing two symbols one on each side of the sign ``='' (read ``equals''), we have a relation called the relation of equality, which means that the two symbols represent the same element. The negation of this relation is obtained by writing the same symbols one on each side of the sign ``≠'' (read ``is not equal to'' or ``is different from'').
\item
  Given a set E and a property of a generic element of E, those of the elements of E which have this property form a new set, called a subset of E. Two equivalent properties therefore define the same subset of E, and conversely.
\end{enumerate}

Let A be a subset of E. When x is a generic element of E, the property ``x belongs to A'' (i.e., ``x is an element of A'') is written `` \(x \in A\) ''; the set of elements which have this property is clearly just A.

The negation of this property is written ``\(x \notin A\)'' (read ``x does not belong to A''); the set of elements of E which have this property is called the complement of A and is written CA or E --- A.

\begin{enumerate}
\def\labelenumi{\arabic{enumi}.}
\setcounter{enumi}{7}
\tightlist
\item
  Some properties, for example \(x = x\) , are true for all elements of E. Any two such properties are equivalent, and the subset they define is the set E itself.
\end{enumerate}

On the other hand, some properties, for example \(x \neq x\) , are not true for any element of E. Again, any two such properties are equivalent, and the subset they define is called the empty subset of E, which is denoted by \(\phi\) .

Note that E and \(\phi\) are complements of each other.

\begin{enumerate}
\def\labelenumi{\arabic{enumi}.}
\setcounter{enumi}{8}
\item
  Let \(a\) be a determinate element of \(E\) . Some properties, for example \(x = a\) , are true only for the single element \(a\) . Any two such properties are equivalent; the subset they define is denoted by \(\{a\}\) , and is called the subset consisting of \(a\) alone.
\item
  The set whose elements are all the subsets of a set E is called the set of subsets of E, and is denoted by \(\mathfrak{P}(\mathbf{E})\) . We have \(\phi \in \mathfrak{P}(\mathbf{E})\) , \(\mathbf{E} \in \mathfrak{P}(\mathbf{E})\) , and \(\{x\} \in \mathfrak{P}(\mathbf{E})\) for all \(x \in E\) . If x denotes a generic element of E, and X a generic element of \(\mathfrak{P}(\mathbf{E})\) , the relation `` \(x \in X\) '' between x and X is called the relation of membership.
\item
  Let \(x\) and \(y\) be two elements of \(\mathbf{E}\) and let \(\mathbf{X}\) be a generic element of \(\mathfrak{P}(\mathbf{E})\) . Then the relation of equality `` \(x = y\) '' is equivalent to the relation ``for all \(\mathbf{X}\) such that \(x \in \mathbf{X}\) , we have \(y \in \mathbf{X}\) ''.
\item
  Let X, Y be two subsets of a set E. If the property \(x \in X\) implies the property \(x \in Y\) , in other words if every element of X belongs to Y, then we say that X is contained in Y, or that Y contains X, or that X is a subset of Y. This relation between X and Y is called the relation of inclusion (of X in Y), and is denoted by ``X ⊂ Y'' or ``Y ⊃ X''. Its negation is denoted by ``X ⊄ Y'' or ``Y ⊅ X''.
\end{enumerate}

For all subsets X of E we have \(\phi \subset X\) and \(X \subset E\) . The relation of membership `` \(x \in X\) '' is equivalent to `` \(\{x\} \subset X\) ''.

The relation ``X ⊂ Y and Y ⊂ Z'' implies ``X ⊂ Z''.

The relation ``X ⊂ Y'' does not exclude the possibility of ``X = Y''. The relation ``X ⊂ Y and Y ⊂ X'' is equivalent to ``X = Y''.

\begin{enumerate}
\def\labelenumi{\arabic{enumi}.}
\setcounter{enumi}{12}
\tightlist
\item
  Let X and Y be any two subsets of E. The set of all elements which have the property `` \(x \in X\) or \(x \in Y\) '' is denoted by X ∪ Y and is called the union of X and Y. The set of all elements which have the property `` \(x \in X\) and \(x \in Y\) '' is denoted by X ∩ Y and is called the intersection of X and Y.
\end{enumerate}

The union and intersection of several subsets of E are defined in the same way.

If \(x, y, z\) are three elements of \(\mathbf{E}\) , the union \(\{x\} \cup \{y\} \cup \{z\}\) is denoted by \(\{x, y, z\}\) . Similarly for any number of (individually named) elements.

Let X and Y be two subsets of E. According as \(X \cap Y \neq \phi\) or \(X \cap Y = \phi\) , we say that X and Y intersect or are disjoint.

\begin{enumerate}
\def\labelenumi{\arabic{enumi}.}
\setcounter{enumi}{13}
\tightlist
\item
  In the statements of the following propositions, X, Y, Z denote any subsets of the same set E.
\end{enumerate}

\begin{enumerate}
\def\labelenumi{(\alph{enumi})}
\item
  We have \(\phi = \mathbb{C}\mathbf{E},\quad \mathbf{E} = \mathbb{C}\phi .\)
\item
  For all \(\mathbf{X}\) we have
\end{enumerate}

\begin{enumerate}
\def\labelenumi{(\arabic{enumi})}
\tightlist
\item
\end{enumerate}

\[
\mathbb {C} (\mathbb {C} \mathbf {X}) = \mathbf {X};\tag{2}
\]

\[
\mathbf {X} \cup \mathbf {X} = \mathbf {X}, \quad \mathbf {X} \cap \mathbf {X} = \mathbf {X};\tag{3}
\]

\[
\mathbf {X} \cup (\mathbb {C X}) = \mathbf {E}, \quad \mathbf {X} \cap (\mathbb {C X}) = \varnothing ;\tag{4}
\]

\[
\mathbf {X} \cup \phi = \mathbf {X}, \quad \mathbf {X} \cap \mathbf {E} = \mathbf {X};\tag{5}
\]

\[
\mathbf {X} \cup \mathbf {E} = \mathbf {E}, \quad \mathbf {X} \cap \phi = \phi .
\]

\begin{enumerate}
\def\labelenumi{(\alph{enumi})}
\setcounter{enumi}{2}
\tightlist
\item
  For all \(\mathbf{X},\mathbf{Y}\) we have
\end{enumerate}

\begin{enumerate}
\def\labelenumi{(\arabic{enumi})}
\setcounter{enumi}{5}
\tightlist
\item
\end{enumerate}

\[
\mathbf {X} \cup \mathbf {Y} = \mathbf {Y} \cup \mathbf {X}, \quad \mathbf {X} \cap \mathbf {Y} = \mathbf {Y} \cap \mathbf {X} \quad (\text { commutativity });\tag{7}
\]

\[
\mathbf {X} \subset \mathbf {X} \cup \mathbf {Y}, \quad \mathbf {X} \cap \mathbf {Y} \subset \mathbf {X};\tag{8}
\]

\[
\mathbf {C} (\mathbf {X} \cup \mathbf {Y}) = (\mathbf {C X}) \cap (\mathbf {C Y}), \quad \mathbf {C} (\mathbf {X} \cap \mathbf {Y}) = (\mathbf {C X}) \cup (\mathbf {C Y}).
\]

\begin{enumerate}
\def\labelenumi{(\alph{enumi})}
\setcounter{enumi}{3}
\item
  The relations \(\mathbf{X} \subset \mathbf{Y}\) , \(\complement \mathbf{X} \supseteq \complement \mathbf{Y}\) , \(\mathbf{X} \cup \mathbf{Y} = \mathbf{Y}\) , \(\mathbf{X} \cap \mathbf{Y} = \mathbf{X}\) are equivalent.
\item
  The relations \(\mathbf{X} \cap \mathbf{Y} = \varnothing\) , \(\mathbf{X} \subset \mathbb{C}\mathbf{Y}\) , \(\mathbf{Y} \subset \mathbb{C}\mathbf{X}\) are equivalent.
\item
  The relations \(\mathbf{X} \cup \mathbf{Y} = \mathbf{E}\) , \(\complement \mathbf{X} \subset \mathbf{Y}\) , \(\complement \mathbf{Y} \subset \mathbf{X}\) are equivalent.
\item
  For all X, Y, Z we have
\end{enumerate}

\[
\mathbf {X} \cup (\mathbf {Y} \cup \mathbf {Z}) = (\mathbf {X} \cup \mathbf {Y}) \cup \mathbf {Z} = \mathbf {X} \cup \mathbf {Y} \cup \mathbf {Z}\tag{9}
\]

\[
\mathrm{X} \cap (\mathrm{Y} \cap \mathrm{Z}) = (\mathrm{X} \cap \mathrm{Y}) \cap \mathrm{Z} = \mathrm{X} \cap \mathrm{Y} \cap \mathrm{Z}\tag{10}
\]

\[
\mathrm{X} \cup (\mathrm{Y} \cap \mathrm{Z}) = (\mathrm{X} \cup \mathrm{Y}) \cap (\mathrm{X} \cup \mathrm{Z})
\]

\[
\mathbf {X} \cap (\mathbf {Y} \cup \mathbf {Z}) = (\mathbf {X} \cap \mathbf {Y}) \cup (\mathbf {X} \cap \mathbf {Z}) \Big \}
\]

(distributivity).

\begin{enumerate}
\def\labelenumi{(\alph{enumi})}
\setcounter{enumi}{7}
\item
  The relation ``X ⊂ Y'' implies the relations ``X ∪ Z ⊂ Y ∪ Z'' and ``X ∩ Z ⊂ Y ∩ Z''.
\item
  The relation ``Z ⊂ X and Z ⊂ Y'' is equivalent to ``Z ⊂ X ∩ Y''. The relation ``X ⊂ Z and Y ⊂ Z'' is equivalent to ``X ∪ Y ⊂ Z''.
\end{enumerate}

\begin{enumerate}
\def\labelenumi{\arabic{enumi}.}
\setcounter{enumi}{14}
\tightlist
\item
  From the identities (8) we conclude that if a subset A of E is obtained from other subsets X, Y, Z of E by applying only the operations \(\mathbb{C}\) , \(\cup\) , \(\cap\) (in any order), then the complement \(\mathbb{C}\mathbf{A}\) can be obtained by replacing the subsets X, Y, Z by their respective complements, and the operations \(\cup\) , \(\cap\) by \(\cap\) , \(\cup\) , respectively, while preserving the order of the operations. This is the duality rule. Given an equality A = B between subsets of the above form, consider the equivalent equality \(\mathbb{C}\mathbf{A} = \mathbb{C}\mathbf{B}\) . If we replace \(\mathbb{C}\mathbf{A}\) and \(\mathbb{C}\mathbf{B}\) by the expressions obtained by applying the duality rule, and if then we replace \(\mathbb{C}\mathbf{X}\) , \(\mathbb{C}\mathbf{Y}\) , \(\mathbb{C}\mathbf{Z}\) by X, Y, Z, respectively, and vice versa, we obtain an equality called the dual of A = B. We can do the same for an inclusion relation A ⊂ B, but then we must take care to replace the sign ``⊂'' by ``⊃''.
\end{enumerate}

The identities above which carry the same number are duals of each other.

\begin{enumerate}
\def\labelenumi{\arabic{enumi}.}
\setcounter{enumi}{15}
\tightlist
\item
  In certain questions we have to consider a fixed subset A of a set E. If X is any arbitrary subset of E, the set \(A \cap X\) is called the trace of X on A, and is sometimes denoted by \(X_{A}\) ; it is considered, in this case, as a subset of A. For all subsets X, Y of E we have
\end{enumerate}

\[
\begin{array}{c} (\mathrm{X} \cup \mathrm{Y}) _ {\mathrm{A}} = \mathrm{X} _ {\mathrm{A}} \cup \mathrm{Y} _ {\mathrm{A}}, \quad (\mathrm{X} \cap \mathrm{Y}) _ {\mathrm{A}} = \mathrm{X} _ {\mathrm{A}} \cap \mathrm{Y} _ {\mathrm{A}}, \\ \mathbb {C} _ {\mathrm{A}} \mathrm{X} _ {\mathrm{A}} = (\mathbb {C} _ {\mathrm{E}} \mathrm{X}) _ {\mathrm{A}}, \end{array}
\]

where \(C_{E}X\) denotes the complement of X in E and \(C_{A}X_{A}\) denotes the complement of \(X_{A}\) in A.

If \(\mathcal{E}\) denotes a set of subsets of \(\mathbf{E}\) , then the set \(\mathcal{E}_{\mathbf{A}}\) of the traces on \(\mathbf{A}\) of the sets of \(\mathcal{E}\) is called the trace of \(\mathcal{E}\) on \(\mathbf{A}\) .

